\subsection{Busca Exata via \textit{Branch-and-Bound}: Fundamentação e Otimizações Graduais (\texttt{v0} \texorpdfstring{$\rightarrow$}{->} \texttt{v3})}
\label{sec:bnb_otimizacoes}

Alimentado pelo limite superior inicial $UB_0$ fornecido por uma das estratégias da
Seção~\ref{sec:warm_starts}, o algoritmo exato resolve o TSP por enumeração implícita com poda
(\textit{Branch-and-Bound}) \cite{levitin2012design}. A seguir, formalizamos primeiro a árvore
de estados e a derivação algébrica do limite inferior que viabilizou nossas otimizações, para
então apresentar em profundidade a progressão controlada das quatro versões implementadas.

\subsubsection{Árvore de Busca e Derivação do Limite Inferior Incremental de Grau 2}
\label{sec:bnb_bound}

Como todo circuito hamiltoniano passa por todos os vértices, fixamos sem perda de generalidade
o vértice de partida $v_0 = 0$. Um estado (ou nó) na profundidade $L \in \{1, \dots, N\}$ da
árvore de busca corresponde a um caminho simples parcial
$s = \langle v_0, v_1, \dots, u, k \rangle$ visitando $L$ vértices distintos, cujo custo
acumulado é $\text{cost}(s) = \sum_{i=0}^{L-2} c(v_i, v_{i+1})$. Ao expandir $s$, o algoritmo
ramifica o problema gerando até $m = N - L$ filhos candidatos, um para cada vértice livre
$w \notin s$. Quando a busca atinge uma folha ($L = N$), basta fechar o ciclo somando a aresta
de retorno $(k, 0)$; se o custo completo for estritamente menor que o incumbente atual $UB$,
atualizamos $UB$ e registramos a nova melhor rota. Para analisar o crescimento dessa árvore,
denotamos por $b$ o \textbf{fator de ramificação efetivo} (número médio de filhos que
sobrevivem à poda por nó expandido, sempre limitado por $N - L$) e por $d \le N$ a profundidade
alcançada pela fronteira ativa de exploração.

\paragraph{Intuição e derivação algébrica do limite inferior por relaxação de grau 2:}
A eficácia do \textit{Branch-and-Bound} repousa na capacidade de estimar, para qualquer caminho
parcial $s$, um \textbf{limite inferior admissível} $LB(s)$ que garanta $LB(s) \le c(H)$ para
todo circuito hamiltoniano $H$ que estenda $s$. Em qualquer circuito válido $H$, cada vértice
$v \in V$ possui exatamente grau $2$ (uma aresta de entrada e uma aresta de saída). Se somarmos
os pesos das duas arestas incidentes a cada vértice ao longo de todos os $N$ vértices, cada
aresta $(x, y) \in H$ será somada exatamente duas vezes (uma na conta de $x$ e outra na conta
de $y$):
\begin{equation}
  c(H) = \frac{1}{2} \sum_{v \in V} \bigl(\text{peso das duas arestas incidentes a } v \text{ em } H\bigr).
\end{equation}
Relaxando a exigência de conectividade global e requerendo apenas que cada vértice escolha duas
arestas incidentes distintas compatíveis com as arestas já fixadas em $s$, obtemos um limite
inferior válido. Para tornar esse cálculo exato e direto, observe que os $N$ vértices do grafo
dividem-se em exatamente três papéis diante de um caminho parcial
$s = \langle 0, v_1, \dots, u, k \rangle$ com $L \ge 2$ vértices:
\begin{enumerate}
  \item \textbf{Vértices internos ($v_1, \dots, u$)}: já possuem ambas as suas arestas no ciclo
  final determinadas pelo caminho $s$ (a aresta que chega e a aresta que sai).
  \item \textbf{Extremidades do caminho ($0$ e $k$)}: cada extremidade já possui exatamente uma
  aresta fixada, a saber $(0, v_1)$ para o vértice $0$ e $(u, k)$ para o vértice $k$. Resta
  escolher a menor \emph{segunda} aresta distinta incidente a cada extremidade. Denotando por
  $\text{min1}[v]$ e $\text{min2}[v]$ os pesos da menor e da segunda menor aresta incidentes a
  $v$, e por $\text{min1\_idx}[v]$ o vizinho que realiza $\text{min1}[v]$, a menor outra aresta
  disponível para uma extremidade $v$ cujo vizinho fixado é $w_{\text{fix}}$ vale exatamente:
  \begin{equation}
    e_{\text{outro}}(v, w_{\text{fix}}) =
    \begin{cases}
      \text{min2}[v], & \text{se } w_{\text{fix}} = \text{min1\_idx}[v], \\
      \text{min1}[v], & \text{caso contrário.}
    \end{cases}
    \label{eq:e_outro}
  \end{equation}
  A intuição é imediata: se a aresta mais curta de $v$ já foi usada na ligação com
  $w_{\text{fix}}$, a segunda aresta distinta terá custo no mínimo $\text{min2}[v]$; caso
  contrário, a menor aresta de $v$ ainda está disponível com peso $\text{min1}[v]$.
  \item \textbf{Vértices ainda livres ($w \notin s$)}: nenhuma de suas duas arestas foi fixada,
  logo sua contribuição mínima possível à soma de graus é $\text{min1}[w] + \text{min2}[w]$.
\end{enumerate}

Ao somar as contribuições desses três grupos, observe que toda aresta fixada $(v_i, v_{i+1})$
do caminho parcial $s$ aparece exatamente duas vezes: uma vez na conta de $v_i$ e outra na
conta de $v_{i+1}$ (sejam eles internos ou extremidades). Assim, a soma de todas as arestas já
fixadas iguala exatamente $2 \cdot \text{cost}(s)$. Reunindo essa parcela às segundas arestas
das extremidades e às duas menores arestas dos vértices livres, obtemos a identidade fechada
para a soma dobrada $S(s)$:
\begin{equation}
  S(s) = 2 \cdot \text{cost}(s) + e_{\text{outro}}(0, v_1) + e_{\text{outro}}(k, u)
         + \underbrace{\sum_{w \notin s} \bigl(\text{min1}[w] + \text{min2}[w]\bigr)}_{R(s)}.
  \label{eq:bound_incremental}
\end{equation}
A consequência algorítmica da Equação~\eqref{eq:bound_incremental} é profunda. Uma formulação
ingênua recalcula o limite inferior varrendo todos os $N$ vértices e suas $N - 1$ arestas a
cada filho gerado, gastando $O(N^2)$ por candidato. Em contraste, pré-computando uma única vez
os vetores $\text{min1}$, $\text{min2}$ e $\text{min1\_idx}$ em $O(N^2 \log N)$ antes da busca
e mantendo apenas o escalar $R(s)$ no estado atual, o passo de estender $s$ até um filho livre
$w$ consiste simplesmente em atualizar $R(s') = R(s) - (\text{min1}[w] + \text{min2}[w])$ e
avaliar a Equação~\eqref{eq:bound_incremental} em tempo estritamente \textbf{$\Theta(1)$}.

\paragraph{Invariante de admissibilidade em métricas contínuas versus inteiras:}
Na apresentação clássica para grafos com pesos inteiros \cite{levitin2012design}, utiliza-se o
arredondamento para cima $LB(s) = \lceil S(s)/2 \rceil$, pois todo circuito completo terá custo
inteiro. Contudo, sob distâncias euclidianas contínuas em ponto flutuante, o operador teto
viola a admissibilidade e pode destruir a otimalidade da busca. Como exemplo concreto, suponha
que o ciclo ótimo verdadeiro de uma instância possua custo $10{,}2$ e que, em um prefixo desse
ótimo, a relaxação forneça $S(s)/2 = 10{,}1 \le 10{,}2$. Aplicar o operador teto produziria
$\lceil 10{,}1 \rceil = 11 > 10{,}2$, podando erroneamente o próprio ramo ótimo caso já se
conheça um incumbente de custo $10{,}8$. Por esse motivo, adotamos estritamente
$LB(s) = S(s)/2$ em instâncias contínuas, reservando o arredondamento ao teto (com margem de
segurança numérica $\epsilon = 10^{-9}$ em $\lceil S(s)/2 - \epsilon \rceil$, para evitar que
um valor exato como $10{,}000000000000002$ em aritmética binária salte para $11$) apenas ao
modo de pesos inteiros.

\subsubsection{Progressão Gradual de Refatorações (\texttt{v0} \texorpdfstring{$\rightarrow$}{->} \texttt{v3})}
\label{sec:bnb_versoes}

Para isolar com clareza o ganho e as limitações de cada intervenção técnica, implementamos
quatro versões incrementais do \textit{Branch-and-Bound}, nas quais cada estágio resolve o
principal gargalo deixado pelo estágio anterior.

\paragraph{Versão \texttt{v0\_naive} (Baseline literal: busca \textit{Best-First} e limite do zero):}
Esta versão reproduz fielmente a formulação clássica de livro-texto \cite{levitin2012design}:
mantém uma fila de prioridades mínima global ordenada pelo valor de $LB(s)$, expandindo sempre
o nó aberto globalmente mais promissor de toda a árvore (estratégia \textit{Best-First}). Cada
estado armazenado na fila contém seu limite inferior, sua profundidade, seu custo acumulado e
uma cópia própria da lista de vértices do caminho parcial. Ao expandir um nó de profundidade
$L$, o algoritmo testa cada vértice $w \in V$ por busca linear na lista do caminho ($O(L)$),
gera uma cópia acrescida da lista ($O(L)$) e recalcula o limite inferior de grau $2$ varrendo
todos os $N$ vértices e suas $N - 1$ arestas ($O(N^2)$). Como há $N - L$ filhos candidatos, o
trabalho por nó expandido atinge $O(N^3)$.

Mais grave que o custo temporal por nó, todavia, é o crescimento espacial da fila global.
Como a estratégia \textit{Best-First} mantém simultaneamente em memória toda a fronteira aberta
de estados promissores enquanto alterna entre subárvores distintas, uma busca que alcance
profundidade $d$ com fator de ramificação efetivo $b > 1$ pode acumular até $O(b^d)$ estados na
fila. Como cada estado carrega sua própria lista de caminho de tamanho $O(N)$, a memória ativa
cresce exponencialmente em $O(b^d \cdot N)$, submetendo o ambiente de execução a forte pressão
do coletor de lixo e ao risco de esgotamento completo da memória principal (\textit{Out of
Memory}, OOM).

\textit{Hipótese $H_0$}: Espera-se que o custo cúbico por expansão e o acúmulo exponencial de
estados na fila \textit{Best-First} levem a versão \texttt{v0\_naive} a tempos impraticáveis ou
ao colapso por falta de memória já nas instâncias pequenas e intermediárias da faixa principal.

\paragraph{Versão \texttt{v1\_gpu} (1ª Otimização: avaliação em lote do limite inferior em GPU):}
A primeira refatoração ataca o gargalo aritmético do recálculo do limite inferior, mantendo
deliberadamente inalterada a estratégia global \textit{Best-First} para permitir uma ablação
limpa do paralelismo de dados. O projeto dessa etapa baseia-se no modelo físico de execução de
unidades gráficas (GPUs). Diferentemente de uma CPU otimizada para baixa latência em fluxos
sequenciais, uma GPU possui milhares de núcleos paralelos projetados para aplicar a mesma
expressão aritmética simultaneamente sobre vetores longos. Porém, cada transferência de dados
entre a memória principal (CPU) e a memória de vídeo (GPU) pelo barramento PCIe, somada ao
disparo de um \textit{kernel}, acarreta uma latência fixa elevada. Enviar à GPU apenas os
$m < 70$ filhos de um único nó gastaria muito mais tempo em comunicação e sincronização do que
no cálculo aritmético propriamente dito.

Para amortizar essa latência, a versão \texttt{v1\_gpu} pré-computa os vetores $\text{min1}$,
$\text{min2}$ e $\text{min1\_idx}$ uma única vez mantendo-os residentes no dispositivo e modifica
o laço principal para retirar da fila de prioridades um \textbf{lote} (\textit{batch}) de até
$B$ nós pais simultaneamente (por exemplo, $B = 128$ ou $1024$). Todas as transições filhas
desses $B$ pais são empacotadas em vetores unidimensionais contíguos de comprimento agregado
$M = \sum_{p=1}^{B} (N - L_p)$, avaliando a Equação~\eqref{eq:bound_incremental} para milhares
de candidatos em uma única chamada vetorizada paralela (executada em GPU via CUDA quando a
aceleração gráfica está habilitada, ou em arranjos vetorizados na CPU em caso contrário). Além
disso, incorporamos um mecanismo de consistência do incumbente: se, durante a montagem do lote,
um dos estados retirados for uma folha ($L = N$) que melhora $UB$, todos os demais pais já
extraídos para o mesmo lote são imediatamente refiltrados contra o novo $UB$ antes da geração
de filhos.

Apesar do ganho aritmético, duas limitações estruturais emergem diretamente da análise teórica.
Primeiro, retirar $B$ pais de uma só vez introduz \textbf{expansão especulativa em lote}: na
busca sequencial estrita, expandir o primeiro pai poderia descobrir rapidamente uma folha
melhor que reduziria $UB$ e podaria os pais seguintes; ao retirá-los juntos em um mesmo lote,
o algoritmo pode expandir estados adicionais que a versão sequencial evitaria, especialmente
sob inicialização frouxa. Segundo, como a fronteira aberta continua residindo na fila global
\textit{Best-First}, a complexidade espacial ativa permanece exponencial em $O(b^d \cdot N)$.

\textit{Hipótese $H_1$}: Prevê-se redução expressiva no custo temporal por candidato avaliado em
relação a \texttt{v0\_naive}, acompanhada possivelmente de aumento no número de nós expandidos
por especulação em lote sob limites iniciais frouxos e da persistência do risco de esgotamento
de memória (OOM) à medida que $N$ cresce, testando se a aceleração puramente aritmética basta
sem reformulação estrutural de memória.

\paragraph{Versão \texttt{v2\_dfs\_memory} (2ª Otimização: refatoração espacial via \textit{Best-Child-First DFS in-place}):}
Para eliminar de forma definitiva o teto de memória exponencial da busca \textit{Best-First},
substituímos a fila global por uma \textbf{busca em profundidade com ordenação local dos filhos
(\textit{Best-Child-First DFS}) operando \textit{in-place}}. Nessa estratégia, apenas um único
caminho da raiz até o nó corrente encontra-se ativo em qualquer instante. Em vez de duplicar
listas para cada filho gerado, a recursão compartilha um único vetor mutável de caminho e um
único vetor booleano de vértices visitados de tamanho $N$: ao descer para um filho $w$,
escreve-se $w$ na posição $L$ do caminho e marca-se $w$ como visitado; ao retornar da chamada
recursiva (\textit{backtracking}), desfaz-se a marcação em $\Theta(1)$.

Uma busca em profundidade arbitrária poderia perder-se por muito tempo em subárvores ruins antes
de encontrar uma rota boa capaz de apertar $UB$. Para preservar a virtude heurística da busca
pelo melhor primeiro sem pagar seu preço em RAM, cada nó calcula a
Equação~\eqref{eq:bound_incremental} para todos os seus filhos livres, descarta imediatamente
aqueles que já satisfazem $LB \ge UB$ (\textbf{poda na geração}) e ordena os sobreviventes em
ordem crescente de limite inferior, descendo primeiro nos ramos mais promissores. Sempre que
uma folha melhora $UB$, ao retornar aos níveis ancestrais cada irmão remanescente é novamente
verificado antes de ser expandido (\textbf{poda na visita}), descartando subárvores inteiras
sem entrar nelas.

\textit{Análise espacial rigorosa}: Com essa reformulação, a profundidade da pilha de recursão
(equivalente ao número de estados ativos simultâneos) fica estritamente limitada a $N$
($\Theta(N)$), e os vetores compartilhados de caminho e visitados ocupam $\Theta(N)$. Mesmo
contabilizando as listas locais de filhos sobreviventes retidas simultaneamente em todos os
níveis ativos da pilha de recursão, a soma no pior caso é
$\sum_{L=1}^{N} (N - L) = N(N - 1)/2 = O(N^2)$ pequenos pares escalares (para $N = 70$, no
máximo $2.415$ pares, da ordem de poucos quilobytes), rompendo por completo qualquer
dependência exponencial em relação ao tamanho da árvore $b^d$.

\textit{Decisão arquitetural sobre a GPU na fronteira de nível 2}: Durante a auditoria dessa
versão, identificou-se que acionar a GPU no interior de cada nó recursivo da DFS seria
prejudicial, pois um nó individual possui no máximo $m = N - L < 70$ filhos, fazendo a
latência de barramento PCIe dominar o tempo total em milhões de chamadas curtas. A solução
arquitetural consiste em deslocar o paralelismo gráfico para a \textbf{fronteira completa de
nível 2}: fixada a origem $v_0 = 0$, há $N - 1$ escolhas para o primeiro passo $v_1$ e $N - 2$
escolhas para o segundo passo $v_2$, formando exatamente $(N - 1)(N - 2)$ subproblemas
independentes de profundidade $3$ (entre $1.482$ e $4.692$ subárvores para
$N \in [40, 70]$). Antes de iniciar a recursão, essa fronteira inteira é avaliada e ordenada em
um único lote paralelo na GPU, podando de imediato subárvores inviáveis e despachando as
sobreviventes em ordem crescente de limite inferior para a DFS recursiva, cujo laço interno
utiliza operações vetorizadas locais na CPU sem viagens ao barramento gráfico.

\textit{Hipótese $H_2$}: A troca para \textit{Best-Child-First DFS in-place} deve eliminar
integralmente as falhas de esgotamento de memória (OOM) em todas as 15 instâncias, mantendo o
pico de estados ativos estritamente limitado por $N$. Como questão empírica aberta, resta
avaliar se o custo fixo de alocar pequenos vetores intermediários a cada chamada recursiva da
DFS limita a vazão de nós por segundo na CPU.

\paragraph{Versão \texttt{v3\_fast\_bound} (3ª Otimização: escalar $\Theta(1)$, \textit{bitmask} e quebra de simetria interna):}
Com a memória controlada em espaço polinomial baixo, o tempo de execução passa a depender de
dois fatores decisivos: a constante multiplicativa do laço mais interno (executado milhões de
vezes) e a redundância combinatória da árvore de busca. A versão \texttt{v3\_fast\_bound}
ataca simultaneamente ambos os fatores por meio de três refatorações acopladas:

\begin{enumerate}
  \item \textbf{Avaliação escalar nativa $\Theta(1)$ sem alocações temporárias}:
  Mesmo operações vetorizadas rápidas em CPU pagam um custo fixo de despacho e criação de
  pequenos objetos de arranjo quando invocadas sobre apenas $5$ a $40$ elementos por nó. Na
  versão \texttt{v3}, convertemos a matriz de distâncias e os vetores pré-computados uma única
  vez para listas escalares nativas da linguagem antes da recursão. Além disso, observamos pela
  Equação~\eqref{eq:bound_incremental} que o termo $e_{\text{outro}}(0, v_1)$ depende apenas do
  primeiro vértice escolhido após a raiz, permanecendo \textbf{constante em toda a subárvore
  enraizada em $v_1$}. Fatorando esse valor fora da recursão, o cálculo do limite inferior de
  cada filho reduz-se a quatro operações aritméticas escalares elementares sem nenhuma alocação
  dinâmica de vetores.

  \item \textbf{Representação dos vértices livres por máscara de bits (\textit{bitmask}) e
  extração direta via \textit{lowbit}}:
  Em um vetor booleano de tamanho $N$, encontrar os $m = N - L$ vértices ainda livres exige
  varrer todas as $N$ posições a cada nó, desperdiçando trabalho linear em $N$ justamente nos
  níveis mais profundos da árvore (onde restam apenas $2$, $3$ ou $4$ vértices livres, mas onde
  se concentra a imensa maioria dos nós expandidos). Para eliminar essa varredura inútil,
  representamos o conjunto de vértices livres em um único inteiro não negativo $M$ (com precisão
  arbitrária nativa na linguagem, suportando $N > 64$ sem risco de estouro), no qual o bit de
  posição $k$ (valor $2^k$) vale $1$ se e somente se o vértice $k$ ainda não foi visitado.
  Marcar ou desmarcar um vértice resume-se a uma operação \texttt{XOR} bit a bit em $\Theta(1)$
  ($M \oplus 2^k$).

  Para percorrer exclusivamente os bits ativos de $M$ sem testar bits zerados, exploramos a
  propriedade do complemento de dois conhecida como \textbf{\textit{lowbit}}: dado o inteiro
  restante \texttt{rem}, a operação $\texttt{lowbit} = \texttt{rem} \mathbin{\&} (-\texttt{rem})$
  isola exatamente o bit ativo menos significativo. Um exemplo binário concreto torna o
  mecanismo transparente: suponha $N = 6$ e que apenas os vértices $\{1, 4\}$ permaneçam livres,
  logo $\texttt{rem} = (010010)_2 = 18$. Em complemento de dois, negar o inteiro inverte todos os
  bits e soma $1$, resultando em $-\texttt{rem} = (\dots 101110)_2$. O \texttt{AND} bit a bit
  entre ambos zera todos os bits acima do primeiro $1$ à direita e preserva exatamente esse bit:
  $(010010)_2 \mathbin{\&} (\dots 101110)_2 = (000010)_2 = 2^1$, revelando o vértice $1$ em
  $\Theta(1)$ por tabela de consulta direta. Aplicando $\texttt{rem} \leftarrow \texttt{rem}
  \oplus \texttt{lowbit}$, salta-se imediatamente para $(010000)_2 = 2^4$ (vértice $4$) sem
  jamais inspecionar os vértices $0$, $2$, $3$ ou $5$. Assim, a iteração sobre os filhos de um
  nó cai de $\Theta(N)$ para $\Theta(m)$, gerando ganho proporcionalmente maior nos níveis mais
  profundos da árvore.

  \item \textbf{Quebra de simetria reversa em nível interno via máscara de bits}:
  Como a distância euclidiana é simétrica ($c(u, v) = c(v, u)$), todo circuito hamiltoniano
  $(0, v_1, v_2, \dots, v_{N-1}, 0)$ possui custo rigorosamente idêntico ao de seu percurso
  inverso $(0, v_{N-1}, \dots, v_2, v_1, 0)$. Para $N \ge 3$, o primeiro vértice após a origem
  ($v_1$) e o último vértice antes do retorno ($v_{N-1}$) são necessariamente distintos; logo,
  entre cada ciclo e seu reverso, exatamente uma das orientações satisfaz a condição canônica
  $v_1 < v_{N-1}$. Restringir a busca a essa orientação reduz o espaço de soluções factíveis
  exatamente pela metade, de $(N - 1)!$ para $(N - 1)!/2$, sem jamais eliminar o ótimo.

  Uma implementação superficial verificaria $v_1 < v_{N-1}$ apenas ao chegar às folhas, já tendo
  desperdiçado todo o trabalho da subárvore. Utilizando a máscara de bits $M$ dos vértices
  livres, conseguimos detectar a inviabilidade da orientação canônica \textbf{ainda nos níveis
  internos da árvore} em tempo $\Theta(1)$. Inicialmente, no nível $1$, proibimos $v_1 = N - 1$
  (pois não existe nenhum vértice maior que $N - 1$ para encerrar o ciclo). Em qualquer nó
  intermediário da subárvore de $v_1$ com $M > 0$, observe a soma geométrica dos pesos de todos
  os bits estritamente menores que $v_1$:
  \begin{equation}
    \sum_{w = 0}^{v_1 - 1} 2^w = 2^{v_1} - 1 < 2^{v_1}.
  \end{equation}
  Em contrapartida, se existir ao menos um vértice ainda livre $w > v_1$, seu bit sozinho já
  contribui com $2^w \ge 2^{v_1 + 1} > 2^{v_1}$, forçando $M > 2^{v_1}$. Portanto, a simples
  comparação inteira em $\Theta(1)$:
  \begin{equation}
    0 < M < 2^{v_1}
    \label{eq:symmetry_bitmask}
  \end{equation}
  é verdadeira \textbf{se e somente se todos os vértices ainda não visitados forem estritamente
  menores que $v_1$}. Para visualizar o ganho em um exemplo concreto, considere $N = 6$ no ramo
  que escolheu $v_1 = 3$ e suponha que, já na profundidade $3$, o caminho tenha visitado os
  vértices $4$ e $5$, restando livres apenas $\{1, 2\}$. A máscara atual vale
  $M = 2^1 + 2^2 = 6$, que é menor que $2^3 = 8$. Qualquer que seja a ordem em que $\{1, 2\}$
  venham a ser visitados, o último vértice $v_5$ será $1$ ou $2$, violando inevitavelmente
  $v_1 < v_5$. Logo, toda a subárvore abaixo desse nó intermediário é podada instantaneamente
  na profundidade $3$, antes mesmo de ser expandida.

  \textit{Prova de corretude}: Seja $H^* = (0, v_1^*, \dots, v_{N-1}^*, 0)$ a orientação
  canônica de um circuito ótimo, satisfazendo $v_1^* < v_{N-1}^*$. Em todo prefixo próprio de
  $H^*$ (níveis $L < N$), o último vértice $v_{N-1}^*$ ainda não foi visitado; logo, o bit
  $2^{v_{N-1}^*}$ permanece ativo em $M$. Como $v_{N-1}^* > v_1^*$, tem-se
  $M \ge 2^{v_{N-1}^*} > 2^{v_1^*}$ ao longo de todo ancestral de $H^*$, de modo que a condição
  \eqref{eq:symmetry_bitmask} jamais poda a orientação canônica do ótimo.
\end{enumerate}

Cabe registrar uma sutileza dinâmica relevante dessa última otimização. Quando o
\textit{Branch-and-Bound} parte de um limite superior inicial frouxo (como no \textit{Cold
Start}), a ordenação pelo limite inferior pode inicialmente conduzir a busca justamente para a
orientação reversa proibida de um ótimo local ou global. Nesse caso, podar essa orientação pode
atrasar momentaneamente a descoberta de uma boa solução incumbente em comparação a uma DFS sem
quebra de simetria, reduzindo temporariamente a poda nos primeiros ramos visitados. Quando a
busca já inicia com um \textit{warm-start} apertado ($UB_0$ próximo de $OPT$), a dependência na
descoberta precoce do incumbente desaparece e o corte de simetria atua diretamente em benefício
da redução de estados.

\textit{Hipótese $H_3$}: Espera-se que \texttt{v3\_fast\_bound} apresente a maior vazão de
candidatos avaliados por segundo entre todas as versões, preservando espaço polinomial e
cortando até metade do espaço de busca sob inicialização apertada, enquanto se investiga
empiricamente se o efeito transitório de atraso na descoberta do incumbente se manifesta sob
\textit{Cold Start}.

A Tabela~\ref{tab:bnb_design} consolida a evolução arquitetural e assintótica das quatro
versões, onde $N$ denota o número de vértices, $m = N - L$ o número de vértices livres na
profundidade $L$, $b$ o fator de ramificação efetivo, $d \le N$ a profundidade máxima da
fronteira aberta e $B$ o tamanho do lote enviado ao dispositivo paralelo.

\begin{table}[htbp]
  \centering
  \caption{Evolução arquitetural e expectativas assintóticas das quatro versões graduais do \textit{Branch-and-Bound}.}
  \label{tab:bnb_design}
  \scriptsize
  \resizebox{\linewidth}{!}{%
  \begin{tabular}{lllllll}
    \toprule
    \textbf{Versão} & \textbf{Busca na Árvore} & \textbf{Bound / Filho} & \textbf{Iteração de Livres} & \textbf{Espaço Ativo} & \textbf{Simetria} & \textbf{Uso da GPU} \\
    \midrule
    \texttt{v0\_naive}      & \textit{Best-First} (fila global) & $O(N^2)$ do zero & Lista $\Theta(N)$ & $O(b^d \cdot N)$ & Não & Não \\
    \texttt{v1\_gpu}        & \textit{Best-First} em lotes ($B$) & $\Theta(1)$ vetor. & Conjunto $\Theta(N)$ & $O(b^d \cdot N)$ & Não & Lotes de $B$ pais \\
    \texttt{v2\_dfs\_memory}& \textit{Best-Child-First DFS} & $\Theta(1)$ vetor. & Vetor \texttt{bool} $\Theta(N)$ & $\Theta(N)$ pilha; $O(N^2)$ local & Não & Fronteira nível 2 \\
    \texttt{v3\_fast\_bound}& \textit{Best-Child-First DFS} & $\Theta(1)$ escalar & \textit{Bitmask} $\Theta(m)$ & $\Theta(N)$ pilha; $O(N^2)$ local & Interna $\Theta(1)$ & Fronteira nível 2 \\
    \bottomrule
  \end{tabular}}
\end{table}
